\documentclass[a4paper, serif, 10pt]{moderncv}

%% moderncv
% style and color
\moderncvstyle{banking}
\moderncvcolor{black}

%% page layout/margins
\usepackage[top=2.5cm, bottom=2.5cm, left=1.5cm, right=1.5cm]{geometry}

\nopagenumbers{}

%% Babel config for Brazilian Portuguese language
% ref:
% - https://latex3.github.io/babel/guides/locale-portuguese.html
\usepackage[brazilian]{babel}

%% fontspec
\usepackage{fontspec}

\IfFontExistsTF{Linux Libertine}{
  \setmainfont[Ligatures={TeX, Common}, Numbers=OldStyle]{Linux Libertine}
}{}
\IfFontExistsTF{Linux Libertine O}{
  \setmainfont[Ligatures={TeX, Common}, Numbers=OldStyle]{Linux Libertine O}
}{}

%% CTeX
% CJK support, used to show CJK characters in the resume
%
% - fontset=none: disable builtin fontset but instead set the CJK font manually
% - heading=false: disable ctex heading
% - punct=kaiming: use kaiming punctuations styles for CJK
% - scheme=plain: use plain scheme, do not override \normalsize font size
% - space=auto: space settings for CJK characters
%
% ref:
% - http://ctan.mirrorcatalogs.com/language/chinese/ctex/ctex.pdf
\usepackage[UTF8, heading=false, punct=kaiming, scheme=plain, space=auto]{ctex}

\IfFontExistsTF{Noto Serif CJK SC}{
  \setCJKmainfont{Noto Serif CJK SC}
}{}
\IfFontExistsTF{Noto Sans CJK SC}{
  \setCJKsansfont{Noto Sans CJK SC}
}{}

%% line spacing
\usepackage{setspace}
\setstretch{1.125}

%% URL styling - use normal text instead of monospace
\urlstyle{same}

% Auto-underline all links
% Defer until \begin{document} so hyperref is already loaded
\AtBeginDocument{%
  \let\oldhref\href
  \renewcommand{\href}[2]{\oldhref{#1}{\underline{#2}}}%
}

%% Patch \httplink and \httpslink to handle full URLs
% The original moderncv macros always prepend http:// or https:// to URLs.
% This causes issues when the URL already has a protocol (e.g., https://example.com)
% resulting in malformed URLs like https://https://example.com.
% This patch checks if the URL already contains "://" and uses it directly if so.
% Note: We use \str_set:Nx (x-expansion) to expand macros like \@homepage first.
\ExplSyntaxOn
\str_new:N \g_yamlresume_url_str

\RenewDocumentCommand{\httpslink}{O{}m}{
  \str_set:Nx \g_yamlresume_url_str {#2}
  \str_if_in:NnTF \g_yamlresume_url_str {://}
    { \url{#2} }
    { \url{https://#2} }
}

\RenewDocumentCommand{\httplink}{O{}m}{
  \str_set:Nx \g_yamlresume_url_str {#2}
  \str_if_in:NnTF \g_yamlresume_url_str {://}
    { \url{#2} }
    { \url{http://#2} }
}
\ExplSyntaxOff

\name{Mariana Souza}{}
\title{Gerente de Produto Sênior | Produtos Digitais, Estratégia e Inovação}
\phone[mobile]{+55 11 98765-4321}
\email{mariana.souza@email.com.br}
\homepage{https://www.linkedin.com/in/mariana-souza-pm}

\address{Av. Paulista, 1000, São Paulo, SP, Brasil, 01310-100}{}{}

\extrainfo{{\small \faLinkedin}\ \href{https://www.linkedin.com/in/mariana-souza-pm}{@mariana-souza-pm} {} {} {} • {} {} {}
{\small \faMedium}\ \href{https://medium.com/@marianasouza}{@marianasouza}}

\begin{document}

\maketitle

\section{Informações Básicas}

\cvline{}{\begin{itemize}
\item Product Manager sênior com 8+ anos de experiência em produtos digitais, atuando em fintechs e plataformas B2B
\item Histórico sólido em liderar ciclos completos de discovery e delivery: da visão e estratégia à entrega ágil e mensuração de resultados
\item Especialista em definição de roadmap orientado a valor, OKRs, priorização com dados e experimentação (A/B testing)
\item Experiência em gestão de times multidisciplinares e interlocução com stakeholders técnicos, comerciais e executivos
\item Apaixonada por resolver problemas reais dos usuários com soluções simples, escaláveis e de alto impacto
\end{itemize}}
\section{Formação}

\cventry{mar. de 2019–dez. de 2020}
        {Mestrado, Gestão de Produtos e Tecnologia}
        {Fundação Getulio Vargas}
        {\url{https://eecsp.fgv.br/}}
        {}
        {\begin{itemize}
\item MBA focado em estratégia de produtos digitais, metodologias ágeis e análise de negócios
\item Desenvolvimento de projetos aplicados a cenários reais de mercado
\end{itemize}
\textbf{Cursos}: Product Strategy, Data-driven Decision Making, Digital Business Models, UX Strategy and Design Thinking}

\cventry{fev. de 2011–dez. de 2015}
        {Bacharelado, Engenharia de Produção, Nota: 8.5}
        {Universidade de São Paulo}
        {\url{https://www.usp.br/}}
        {}
        {\begin{itemize}
\item Formação com ênfase em otimização de processos, gestão de projetos e análise de dados
\item Iniciação científica em modelagem de processos aplicada à indústria de serviços
\end{itemize}
\textbf{Cursos}: Gestão de Processos, Pesquisa Operacional, Estatística Aplicada, Economia e Finanças}
\section{Experiência Profissional}

\cventry{jan. de 2021–Atual}
        {Gerente de Produto Sênior}
        {Nexo Bank}
        {\url{https://www.nexobank.com.br}}
        {}
        {\begin{itemize}
\item Responsável pela estratégia e roadmap do produto de conta digital para pessoa física, com mais de 3 milhões de usuários ativos
\item Lidero squads multifuncionais de engenharia, design e dados na entrega de funcionalidades de pagamento, investimento e crédito
\item Implantei ciclo de discovery contínuo com pesquisa qualitativa e quantitativa, elevando a taxa de sucesso das entregas em 35\%
\item Defini OKRs e métricas de produto (activation, retention, NPS) em conjunto com stakeholders executivos e comerciais
\item Coordeno o processo de priorização com base em impacto, esforço e viabilidade técnica, aumentando a velocidade de aprendizado do time
\end{itemize}
\textbf{Palavras-chave}: Fintech, Produtos Bancários, OKRs, Discovery, A/B Testing}

\cventry{mar. de 2018–dez. de 2020}
        {Gerente de Produto}
        {Payline Tecnologia}
        {\url{https://www.payline.com.br}}
        {}
        {\begin{itemize}
\item Estruturei a área de produto de uma plataforma B2B de pagamentos recorrentes, responsável pela visão e estratégia do produto
\item Conduzi mais de 50 entrevistas com clientes e análises de uso para identificar dores e oportunidades de melhoria
\item Liderei o lançamento de um novo módulo de conciliação que reduziu em 40\% o tempo de fechamento financeiro dos clientes
\item Atuei na formação do time de product designers e analistas de dados, promovendo uma cultura orientada a métricas
\end{itemize}
\textbf{Palavras-chave}: SaaS B2B, Entrevistas com Usuários, Roadmap, Analytics, Liderança de Produto}

\cventry{jan. de 2016–fev. de 2018}
        {Analista de Produto}
        {Loja Ágil}
        {\url{https://www.lojaagil.com.br}}
        {}
        {\begin{itemize}
\item Atuei no product discovery de um e-commerce de moda, apoiando a definição de funcionalidades e priorização do backlog
\item Criei dashboards de acompanhamento de indicadores-chave (conversão, ticket médio, churn) utilizando SQL e ferramentas de BI
\item Colaborei com times de UX e engenharia na implementação de melhorias de interface e jornada de compra
\end{itemize}
\textbf{Palavras-chave}: E-commerce, Backlog, SQL, Indicadores, UX}
\section{Idiomas}

\cvline{Português}{Nativo ou Bilíngue \hfill \textbf{Palavras-chave}: Fala nativa}
\cvline{Inglês}{Proficiência Profissional Fluente \hfill \textbf{Palavras-chave}: TOEFL iBT 105}
\cvline{Espanhol}{Proficiência Profissional Mínima \hfill \textbf{Palavras-chave}: Espanhol conversacional}
\section{Competências}

\cvline{Gestão de Produto}{Especialista \hfill \textbf{Palavras-chave}: Product Discovery, Roadmap, OKRs, Priorização, Product Strategy, User Research, A/B Testing, Métricas de Produto}
\cvline{Métodos Ágeis}{Avançado \hfill \textbf{Palavras-chave}: Scrum, Kanban, Lean, Design Thinking}
\cvline{Dados e Analytics}{Avançado \hfill \textbf{Palavras-chave}: SQL, Amplitude, Mixpanel, Google Analytics, Tableau}
\cvline{Liderança e Comunicação}{Especialista \hfill \textbf{Palavras-chave}: Facilitação, Negociação, Comunicação Executiva, Mentoria}
\section{Prêmios}

\cventry{nov. de 2023}
        {Associação Brasileira de Produtos Digitais}
        {Prêmio Produto do Ano}
        {}
        {}
        {Reconhecimento pela liderança no redesenho do onboarding do banco digital, elevando a ativação de novos usuários em 28\%.}
\section{Certificados}

\cventry{jun. de 2020}
        {Scrum Alliance}
        {Certified Scrum Product Owner (CSPO)}
        {\url{https://www.scrumalliance.org/}}
        {}
        {}

\cventry{mar. de 2022}
        {Amplitude}
        {Product Analytics Certification}
        {\url{https://amplitude.com/}}
        {}
        {}
\section{Publicações}

\cventry{mai. de 2022}
        {Como priorizar funcionalidades com base em dados e empatia}
        {Medium}
        {\url{https://medium.com/@marianasouza/como-priorizar-funcionalidades-com-base-em-dados-e-empatia}}
        {}
        {Artigo sobre a combinação de métricas quantitativas e pesquisas qualitativas para definir o que construir primeiro,
com exemplos práticos de priorização em times de produto no Brasil.}
\section{Projetos}

\cventry{jun. de 2019–nov. de 2020}
        {Redesenho do portal de clientes com foco em autosserviço e redução de chamados de suporte.}
        {Portal de Autosserviço Payline}
        {\url{https://www.payline.com.br/portal}}
        {}
        {\begin{itemize}
\item Discovery e prototipação de novas jornadas de autosserviço
\item Definição de roadmap e coordenação de 3 squads na entrega do novo portal
\item Resultado: redução de 32\% nos chamados de suporte e aumento da satisfação do cliente (CSAT)
\end{itemize}
\textbf{Palavras-chave}: B2B, Autosserviço, UX, Roadmap, CSAT}
\section{Interesses}

\cvline{Produto e Inovação}{Inteligência Artificial, Plataformas Digitais, Experiência do Usuário}
\cvline{Estilo de Vida}{Corrida, Yoga, Culinária}
\section{Voluntariado}

\cventry{jan. de 2022–Atual}
        {Mentora de Product Managers}
        {Mulheres em Produto}
        {\url{https://www.mulheresemproduto.com.br}}
        {}
        {\begin{itemize}
\item Mentoria mensal para profissionais em início de carreira em gestão de produto
\item Apoio na construção de carreira, preparação para entrevistas e desenvolvimento de habilidades de liderança
\end{itemize}}

\end{document}